
\documentclass[11pt]{article}
\usepackage[margin=1in]{geometry}
\usepackage[T1]{fontenc}
\usepackage[utf8]{inputenc}
\usepackage{lmodern}
\usepackage{amsmath}
\usepackage{amssymb}
\usepackage{graphicx}
\usepackage{booktabs}
\usepackage{float}
\usepackage{hyperref}
\hypersetup{hidelinks}
\setlength{\parskip}{0.5em}
\setlength{\parindent}{0pt}
\begin{document}
\sloppy
\section{Step 7: Duration-Aware, Lagged ROI Analysis with Variable Trial Windows}
\label{sec:step7_variable_window}

\subsection{Objective}
The objective of Step~7 is to replace the fixed post-onset summary window used in Steps~5 and~6
with a \textbf{question-specific variable analysis window} that depends on the actual duration of each trial.

The goal of this step is to test whether the abstract-versus-concrete effect in the left frontal ROI
remains when the fNIRS summary is computed over a trial window that:
\begin{itemize}
    \item starts from the actual question onset;
    \item ends at the actual participant answer;
    \item includes a haemodynamic lag;
    \item varies from trial to trial instead of being fixed at the same length for every question.
\end{itemize}

This step should remain simple.
It should produce:
\begin{itemize}
    \item one primary covariate-adjusted inferential model;
    \item exact p-values;
    \item confidence intervals;
    \item effect estimates;
    \item a short final conclusion.
\end{itemize}

\subsection{Why Step~7 Is Needed}
Step~6 showed that the fixed-window front ROI analysis was close to, but still above, the conventional
significance threshold.
At the same time, the descriptive visualizations suggested that:
\begin{itemize}
    \item condition differences evolve over time rather than appearing in one perfectly fixed segment;
    \item the waveform shape varies across channels;
    \item the questions differ in length, structure, and difficulty.
\end{itemize}

Therefore, a variable-window analysis is a natural next step because it allows the fNIRS summary
to follow the real duration of the question-processing period for each trial.

\subsection{Core Analytical Decision}
To keep Step~7 interpretable, the following elements should remain fixed from Step~6:
\begin{enumerate}
    \item the same preprocessed continuous fNIRS files from Step~3;
    \item the same included participant-session cohort carried forward from Step~6;
    \item the same long-channel-only rule;
    \item the same frontal ROI as the primary ROI;
    \item the same short-channel exclusion from inference;
    \item the same question-level covariates;
    \item the same mixed-effects modeling framework with participant and question random intercepts.
\end{enumerate}

The main change in Step~7 is:
\begin{quote}
the trial outcome is no longer based on a fixed 7--11~s window, but on a lagged
question-specific interval that varies with the actual trial duration.
\end{quote}

\subsection{Primary Scientific Question}
The primary Step~7 question is:
\begin{quote}
In the left frontal ROI, does the abstract-versus-concrete effect remain after adjusting for
question length and item difficulty when the fNIRS response is summarized over a duration-aware,
lagged, trial-specific interval rather than a fixed window?
\end{quote}

\subsection{Why Continuous Preprocessed Data Must Be Used}
Step~5 and Step~6 were based on onset-locked epochs extending from:
\begin{verbatim}
-3.5 s to 13.2 s around question onset
\end{verbatim}

That epoch structure was appropriate for fixed-window analyses.
However, Step~7 should not rely on those fixed-length epochs as the main analysis input,
because a variable trial window may extend beyond the end of that epoch for slower responses.

Therefore, Step~7 should use:
\begin{quote}
the continuous preprocessed raw data from Step~3 together with the harmonized trigger table.
\end{quote}

\subsection{Primary ROI Definition}
The primary ROI remains the same left frontal ROI used in Steps~5 and~6.

\subsubsection{Front ROI HbO channels}
\begin{verbatim}
S2_D6 hbo
S2_D4 hbo
S1_D6 hbo
S1_D3 hbo
S5_D6 hbo
S5_D4 hbo
S5_D3 hbo
\end{verbatim}

\subsubsection{Back ROI HbO channels}
The back ROI is retained as a secondary ROI:
\begin{verbatim}
S4_D2 hbo
S3_D1 hbo
S6_D2 hbo
S6_D1 hbo
S6_D5 hbo
S7_D2 hbo
S7_D5 hbo
S8_D1 hbo
S8_D5 hbo
\end{verbatim}

\subsubsection{Transition channels excluded from the primary confirmatory summaries}
\begin{verbatim}
S5_D7 hbo
S4_D4 hbo
S4_D7 hbo
S3_D3 hbo
S3_D7 hbo
S6_D7 hbo
\end{verbatim}

\subsubsection{Short-separation pairs excluded from inference}
\begin{verbatim}
S1_D8
S2_D9
S3_D10
S4_D11
S5_D12
S6_D13
S7_D14
S8_D15
\end{verbatim}

\subsection{Primary Trial Window Definition}
For participant \(i\) and trial \(k\), let:
\begin{itemize}
    \item \(t^{\mathrm{on}}_{ik}\) denote the harmonized question onset time;
    \item \(t^{\mathrm{ans}}_{ik}\) denote the harmonized answer time;
    \item \(t^{\mathrm{next}}_{ik}\) denote the next question onset time, if it exists.
\end{itemize}

The raw behavioral duration of the trial is:
\begin{equation}
L_{ik} = t^{\mathrm{ans}}_{ik} - t^{\mathrm{on}}_{ik}.
\end{equation}

The primary Step~7 analysis window should be a \textbf{lagged duration-aware interval}:
\begin{equation}
W^{(\delta)}_{ik}
=
\left[
t^{\mathrm{on}}_{ik} + \delta,\;
\min\!\bigl(t^{\mathrm{ans}}_{ik} + \delta,\; t^{\mathrm{next}}_{ik} - 0.5\bigr)
\right],
\end{equation}
where the primary lag is:
\begin{equation}
\delta = 4 \text{ s}.
\end{equation}

If the trial is the last question before the end of a session, then \(t^{\mathrm{next}}_{ik}\) should be
replaced by the end of the available continuous recording.

\subsection{Why This Window Is Better Than a Fixed One}
This window is preferable for Step~7 because:
\begin{itemize}
    \item it respects the actual duration of the question-processing period for each trial;
    \item it incorporates a simple haemodynamic lag;
    \item it avoids contamination from the next trial by truncating at the next onset when necessary;
    \item it changes the interval length with the trial instead of forcing every trial into the same fixed window.
\end{itemize}

\subsection{Baseline Definition}
For each trial, define the baseline interval as:
\begin{equation}
B_{ik} = [t^{\mathrm{on}}_{ik} - 2,\; t^{\mathrm{on}}_{ik}].
\end{equation}

If the full two-second baseline is unavailable, the trial should be excluded from the primary analysis
and logged.

\subsection{Primary Trial-Level Outcome}
Let \(x_{ijk}(t)\) denote the preprocessed HbO signal for participant \(i\), channel \(j\), and trial \(k\).

The baseline-corrected trial response for a channel is:
\begin{equation}
r_{ijk}
=
\overline{x_{ijk}(t \in W^{(\delta)}_{ik})}
-
\overline{x_{ijk}(t \in B_{ik})}.
\end{equation}

The critical point is that this is the \textbf{mean} signal over the variable window, not the raw area.
This avoids making longer trials mechanically larger simply because they last longer.

\subsection{Front ROI Trial Summary}
For each trial, the front ROI response is:
\begin{equation}
y^{(F)}_{ik}
=
\frac{1}{J^{(F)}_{ik}}
\sum_{j \in F_{ik}} r_{ijk},
\end{equation}
where \(F_{ik}\) is the set of valid front-ROI channels available for participant \(i\) on trial \(k\),
and \(J^{(F)}_{ik}\) is the number of such channels.

\subsection{Back ROI Trial Summary}
For secondary analyses, define:
\begin{equation}
y^{(B)}_{ik}
=
\frac{1}{J^{(B)}_{ik}}
\sum_{j \in B_{ik}} r_{ijk},
\end{equation}
where \(B_{ik}\) is the set of valid back-ROI channels available for participant \(i\) on trial \(k\).

\subsection{Primary Covariates}
Step~6 showed that the structural language variables were not strongly imbalanced by condition,
but ENEM item correctness differed more substantially across abstract and concrete questions.
Therefore, Step~7 should treat both question size and question difficulty explicitly.

The primary covariates should be:
\begin{itemize}
    \item standardized log total word count;
    \item standardized ENEM item correctness.
\end{itemize}

Define:
\begin{equation}
w_k = z\!\left(\log(\texttt{total\_word\_count}_k)\right)
\end{equation}
and
\begin{equation}
c_k = z(\texttt{correctness}_k).
\end{equation}

\subsection{Primary Model}
Let:
\begin{equation}
A_{ik} =
\begin{cases}
1 & \text{if the trial is Abstract} \\
0 & \text{if the trial is Concrete}
\end{cases}
\end{equation}

The primary Step~7 model should be:
\begin{equation}
y^{(F)}_{ik}
=
\beta_0
+
\beta_1 A_{ik}
+
\beta_2 w_k
+
\beta_3 c_k
+
u_i
+
v_k
+
\varepsilon_{ik},
\end{equation}
where:
\begin{itemize}
    \item \(u_i \sim \mathcal{N}(0,\sigma_u^2)\) is the participant random intercept;
    \item \(v_k \sim \mathcal{N}(0,\sigma_v^2)\) is the question random intercept;
    \item \(\varepsilon_{ik} \sim \mathcal{N}(0,\sigma^2)\) is the residual term.
\end{itemize}

\subsection{Primary Inferential Target}
The primary inferential target is:
\begin{equation}
\beta_1,
\end{equation}
which represents the adjusted abstract-versus-concrete effect in the left frontal ROI under the
variable-window definition.

\subsection{Primary Hypothesis}
The primary hypotheses are:

\paragraph{Null hypothesis}
\begin{equation}
H_0: \beta_1 = 0
\end{equation}

\paragraph{Alternative hypothesis}
\begin{equation}
H_1: \beta_1 \neq 0
\end{equation}

The significance threshold should remain:
\begin{verbatim}
alpha = 0.05
\end{verbatim}

\subsection{Trial Inclusion Rules}
A trial should be included in the primary Step~7 analysis if:
\begin{itemize}
    \item it belongs to a participant-session already included in Step~6;
    \item it has a valid harmonized question onset;
    \item it has a valid harmonized answer time;
    \item the behavioral duration \(L_{ik}\) is positive;
    \item the two-second baseline interval is fully available;
    \item the final lagged variable window is fully available;
    \item the final lagged variable window length is at least 2.0~s after any truncation;
    \item at least one valid front ROI channel is present.
\end{itemize}

\subsection{Participant Inclusion Rules}
A participant-session should be retained in the primary Step~7 analysis if:
\begin{itemize}
    \item it was included in Step~6;
    \item it contributes at least two valid abstract trials;
    \item it contributes at least two valid concrete trials;
    \item at least one valid front ROI channel remains for the included trials.
\end{itemize}

\subsection{Pre-Model Duration Diagnostics}
Before fitting the mixed models, Step~7 should include a descriptive analysis of how trial duration
relates to condition and question properties.

At minimum, this pre-model diagnostic should report:
\begin{itemize}
    \item response-time distribution by condition;
    \item correlation of response time with total word count;
    \item correlation of response time with sentence count;
    \item correlation of response time with sentence length;
    \item correlation of response time with ENEM item correctness.
\end{itemize}

This is useful because it quantifies the very mechanism that motivates the variable-window analysis.

\subsection{Alternative Structure Models}
To assess whether the result depends on the choice of structural covariate, Step~7 should also fit:

\subsubsection{Alternative Model A: Sentence-count model}
\begin{equation}
y^{(F)}_{ik}
=
\beta_0
+
\beta_1 A_{ik}
+
\beta_2 s_k
+
\beta_3 c_k
+
u_i
+
v_k
+
\varepsilon_{ik},
\end{equation}
where
\begin{equation}
s_k = z(\texttt{sentence\_count}_k).
\end{equation}

\subsubsection{Alternative Model B: Sentence-length model}
\begin{equation}
y^{(F)}_{ik}
=
\beta_0
+
\beta_1 A_{ik}
+
\beta_2 \ell_k
+
\beta_3 c_k
+
u_i
+
v_k
+
\varepsilon_{ik},
\end{equation}
where
\begin{equation}
\ell_k = z(\texttt{sentence\_length}_k).
\end{equation}

\subsection{Secondary Back ROI Model}
The same primary variable-window model should be fit for the back ROI as a secondary analysis:
\begin{equation}
y^{(B)}_{ik}
=
\beta_0
+
\beta_1 A_{ik}
+
\beta_2 w_k
+
\beta_3 c_k
+
u_i
+
v_k
+
\varepsilon_{ik}.
\end{equation}

\subsection{Secondary Dissociation Model}
To test whether the condition effect is more frontal than posterior under the variable-window definition,
Step~7 should define a trial-level dissociation outcome:
\begin{equation}
d_{ik} = y^{(F)}_{ik} - y^{(B)}_{ik}.
\end{equation}

The dissociation model should then be:
\begin{equation}
d_{ik}
=
\beta_0
+
\beta_1 A_{ik}
+
\beta_2 w_k
+
\beta_3 c_k
+
u_i
+
v_k
+
\varepsilon_{ik}.
\end{equation}

This is a secondary analysis.

\subsection{Sensitivity Analyses}
To keep Step~7 robust but still simple, the following sensitivity analyses should be included.

\subsubsection{Sensitivity 1: Lag choice}
Repeat the primary front ROI model using:
\begin{itemize}
    \item \(\delta = 3\) s;
    \item \(\delta = 4\) s;
    \item \(\delta = 5\) s.
\end{itemize}

\subsubsection{Sensitivity 2: Cap long windows}
Repeat the primary model using a capped lagged window:
\begin{equation}
W^{(\delta,\mathrm{cap})}_{ik}
=
\left[
t^{\mathrm{on}}_{ik} + \delta,\;
\min\!\bigl(t^{\mathrm{ans}}_{ik} + \delta,\; t^{\mathrm{next}}_{ik} - 0.5,\; t^{\mathrm{on}}_{ik} + 16\bigr)
\right].
\end{equation}

This checks whether very long trials disproportionately influence the result.

\subsubsection{Sensitivity 3: Response-time adjustment}
Because response time is closely tied to the window length, it should not be included in the primary model.
However, one sensitivity model may include standardized log response time:
\begin{equation}
r_{ik} = z\!\left(\log(L_{ik})\right)
\end{equation}
and
\begin{equation}
y^{(F)}_{ik}
=
\beta_0
+
\beta_1 A_{ik}
+
\beta_2 w_k
+
\beta_3 c_k
+
\beta_4 r_{ik}
+
u_i
+
v_k
+
\varepsilon_{ik}.
\end{equation}

This model should be interpreted cautiously because response time may be partly downstream of condition.

\subsubsection{Sensitivity 4: Fixed-window benchmark}
Repeat the Step~6 primary front ROI model alongside the variable-window model so that the two
approaches can be compared directly in the same report.

\subsection{Step-by-Step Workflow}

\subsubsection{Step~7.1: Freeze the Step~6 cohort}
Reuse the same included participant-session set as Step~6 unless a documented file or QC problem
requires explicit re-review.

\subsubsection{Step~7.2: Return to the continuous preprocessed data}
Do not use the fixed-length Step~5/6 epochs as the primary input.
Load the continuous preprocessed data from Step~3 together with the harmonized trigger table.

\subsubsection{Step~7.3: Build the trial-level timing table}
For each valid trial, store:
\begin{itemize}
    \item participant ID;
    \item question ID;
    \item question condition;
    \item question onset time;
    \item answer time;
    \item next question onset time;
    \item raw answer latency;
    \item final lagged variable-window start;
    \item final lagged variable-window end;
    \item final variable-window duration.
    \item any truncation flag.
\end{itemize}

\subsubsection{Step~7.4: Compute baseline-corrected ROI trial summaries}
For each included trial:
\begin{itemize}
    \item compute the front ROI mean over the variable window;
    \item compute the back ROI mean over the variable window;
    \item record the number of valid channels contributing to each ROI.
\end{itemize}

\subsubsection{Step~7.5: Merge question covariates}
Merge:
\begin{itemize}
    \item total word count;
    \item sentence count;
    \item sentence length;
    \item ENEM item correctness.
\end{itemize}

\subsubsection{Step~7.6: Run duration diagnostics}
Produce descriptive plots and tables showing how response time varies by condition and by question properties.

\subsubsection{Step~7.7: Fit the primary front ROI variable-window model}
Fit the primary mixed-effects model:
\begin{itemize}
    \item condition;
    \item log word count;
    \item ENEM correctness;
    \item random intercept for participant;
    \item random intercept for question.
\end{itemize}

\subsubsection{Step~7.8: Fit the alternative structure models}
Fit the sentence-count and sentence-length alternatives.

\subsubsection{Step~7.9: Fit the secondary back and dissociation models}
Fit the variable-window back ROI model and the variable-window dissociation model.

\subsubsection{Step~7.10: Run the sensitivity analyses}
Run:
\begin{itemize}
    \item lag sensitivity;
    \item capped-window sensitivity;
    \item response-time sensitivity;
    \item fixed-window benchmark comparison.
\end{itemize}

\subsection{Primary Reporting Quantities}
The primary Step~7 results table must report:
\begin{itemize}
    \item number of included participants;
    \item number of included questions;
    \item number of included trials;
    \item mean and distribution of variable-window durations;
    \item adjusted condition coefficient \(\beta_1\);
    \item standard error;
    \item test statistic;
    \item p-value;
    \item 95\% confidence interval;
    \item word-count coefficient;
    \item ENEM-correctness coefficient;
    \item participant random-intercept variance;
    \item question random-intercept variance.
\end{itemize}

\subsection{Secondary Reporting Quantities}
The secondary Step~7 tables should report:
\begin{itemize}
    \item back ROI variable-window model results;
    \item dissociation model results;
    \item sentence-count alternative model results;
    \item sentence-length alternative model results;
    \item lag sensitivity results;
    \item capped-window sensitivity results;
    \item response-time sensitivity results;
    \item fixed-window benchmark comparison.
\end{itemize}

\subsection{Primary Conclusion Rule}
The final Step~7 conclusion should be based first on the condition coefficient \(\beta_1\)
from the primary variable-window front ROI model.

The interpretation rule should be:
\begin{itemize}
    \item if the adjusted condition coefficient is significant and positive, conclude that abstract
    questions are associated with a larger frontal HbO response than concrete questions even when
    the response is summarized over a lagged question-specific duration-aware interval;
    \item if the adjusted condition coefficient is significant and negative, conclude that concrete
    questions are associated with a larger frontal HbO response than abstract questions under the
    variable-window definition;
    \item if the adjusted condition coefficient is not significant, conclude that the variable-window
    adjusted model does not show statistically significant evidence for a frontal abstract-versus-concrete effect.
\end{itemize}

\subsection{Conclusion Templates}

\paragraph{Template if the primary effect is significant and abstract \(>\) concrete.}
In the duration-aware mixed-effects analysis, abstract questions were associated with a significantly
larger left frontal ROI HbO response than concrete questions after adjusting for log word count and
ENEM item correctness
(\(\beta=\ldots\), \(SE=\ldots\), \(p=\ldots\), 95\% CI \([\ldots,\ldots]\)).
This indicates that the frontal condition effect remains when the analysis window is allowed to vary
with the actual duration of each question.

\paragraph{Template if the primary effect is significant and concrete \(>\) abstract.}
In the duration-aware mixed-effects analysis, concrete questions were associated with a significantly
larger left frontal ROI HbO response than abstract questions after adjusting for log word count and
ENEM item correctness
(\(\beta=\ldots\), \(SE=\ldots\), \(p=\ldots\), 95\% CI \([\ldots,\ldots]\)).

\paragraph{Template if the primary effect is not significant.}
In the duration-aware mixed-effects analysis, the abstract-versus-concrete effect in the left frontal ROI
was not statistically significant after adjusting for log word count and ENEM item correctness
(\(\beta=\ldots\), \(SE=\ldots\), \(p=\ldots\), 95\% CI \([\ldots,\ldots]\)).
This suggests that allowing the trial window to vary with question duration does not by itself yield
statistically significant evidence for a frontal condition effect in the current dataset.

\subsection{Expected Outputs}

\subsubsection*{Tables}
\begin{verbatim}
01_step7_trial_timing_table.csv
02_step7_duration_diagnostics.csv
03_step7_trial_variable_window_summary.csv
04_step7_primary_front_variable_window_model.csv
05_step7_front_sentencecount_model.csv
06_step7_front_sentencelength_model.csv
07_step7_back_variable_window_model.csv
08_step7_dissociation_variable_window_model.csv
09_step7_lag_sensitivity.csv
10_step7_capped_window_sensitivity.csv
11_step7_rt_sensitivity.csv
12_step7_fixed_vs_variable_benchmark.csv
13_step7_exclusion_log.csv
14_step7_final_conclusion.csv
\end{verbatim}

\subsubsection*{Figures}
\begin{verbatim}
figures/step7/response_time_by_condition.png
figures/step7/response_time_vs_wordcount.png
figures/step7/response_time_vs_correctness.png
figures/step7/variable_window_duration_histogram.png
figures/step7/front_roi_partial_effect_variable_window.png
figures/step7/back_roi_partial_effect_variable_window.png
figures/step7/dissociation_partial_effect_variable_window.png
figures/step7/fixed_vs_variable_comparison.png
figures/step7/model_residuals_primary.png
\end{verbatim}

\subsubsection*{Reports}
\begin{verbatim}
reports/step7/step7_variable_window_report.tex
reports/step7/tables/primary_model_table.tex
reports/step7/tables/secondary_model_table.tex
reports/step7/text/final_conclusion.tex
\end{verbatim}

\subsection{Minimal Figures for the Main Report}
The main Step~7 report should include at least:
\begin{enumerate}
    \item response-time distribution by condition;
    \item response time versus total word count;
    \item response time versus ENEM item correctness;
    \item histogram of final variable-window durations;
    \item partial-effect plot for the primary front ROI variable-window model;
    \item direct benchmark comparison of the Step~6 fixed-window coefficient and the Step~7 variable-window coefficient.
\end{enumerate}

\subsection{Quality-Control Checks}
Before Step~7 is considered complete, verify that:
\begin{enumerate}
    \item the Step~6 cohort is unchanged unless explicitly re-reviewed;
    \item the continuous preprocessed files are used rather than truncated onset epochs;
    \item every included trial has valid onset and answer times;
    \item every included trial has a fully available baseline interval;
    \item the variable window does not overlap the next trial onset;
    \item the final variable window length is positive and at least 2~s;
    \item only long-separation channels are used in the inferential ROI summaries;
    \item the short-separation pairs remain excluded from inference;
    \item the front ROI channel list is identical to Step~6;
    \item the primary model includes both question size and question difficulty;
    \item the primary model converges;
    \item all sensitivity analyses are clearly labeled as non-primary.
\end{enumerate}

\subsection{Fallback Rule}
If the primary mixed-effects model fails to converge reliably, the fallback analysis should be:
\begin{quote}
a linear model with participant fixed effects and question fixed effects, keeping the same
variable-window front ROI outcome and the same primary covariates.
\end{quote}

\subsection{Minimal Completion Criteria}
Step~7 is complete when:
\begin{itemize}
    \item the variable trial windows have been constructed;
    \item the duration diagnostics have been completed;
    \item the primary front ROI variable-window model has been fit;
    \item the alternative structure models have been fit;
    \item the secondary back and dissociation models have been fit;
    \item the lag and capped-window sensitivities have been completed;
    \item the fixed-window versus variable-window benchmark comparison has been completed;
    \item the primary adjusted p-value has been reported;
    \item a short final conclusion has been generated;
    \item all exclusions and warnings have been documented.
\end{itemize}

\subsection{Short Practical Summary}
In simple terms, Step~7 should:
\begin{quote}
reuse the same cohort and the same frontal ROI from Step~6, move from a fixed 7--11~s summary
to a lagged question-specific variable window that runs with the actual duration of each trial,
adjust for both question length and ENEM difficulty, and test whether the abstract-versus-concrete
effect survives under this more duration-aware definition.
\end{quote}
\end{document}
